% ============================================================
% f2013_shaanxi_arts_q20_prism.tex
% 2013 陕西卷（文史类）第 20 题 四棱柱（TikZ 矢量重绘）
% 底面 ABCD 为正方形（斜二测示意），O 为底面中心，A1O ⊥ 底面
% 虚线：被遮挡的底边、侧棱、两平面的对角线以及 A1O
% ============================================================
\begin{tikzpicture}[
  x=1cm, y=1cm, line width=0.5pt, font=\footnotesize, >=Stealth
]

  % ---------- 底面（正方形示意） ----------
  \draw (0,0) -- (3,0) -- (4.6,1.2);
  \draw[dashed] (4.6,1.2) -- (1.6,1.2) -- (0,0);

  % ---------- 侧棱 ----------
  \draw (0,0) -- (2.3,3.0);         % AA1
  \draw (3,0) -- (5.3,3.0);         % BB1
  \draw (4.6,1.2) -- (6.9,4.2);     % CC1
  \draw[dashed] (1.6,1.2) -- (3.9,4.2);   % DD1

  % ---------- 上底面 ----------
  \draw (2.3,3.0) -- (5.3,3.0) -- (6.9,4.2) -- (3.9,4.2) -- cycle;

  % ---------- 平面 A1BD 与 平面 CD1B1 的对角线 ----------
  \draw[dashed] (2.3,3.0) -- (3,0);
  \draw[dashed] (2.3,3.0) -- (1.6,1.2);
  \draw[dashed] (3,0) -- (1.6,1.2);
  \draw[dashed] (4.6,1.2) -- (3.9,4.2);
  \draw[dashed] (4.6,1.2) -- (5.3,3.0);
  \draw[dashed] (3.9,4.2) -- (5.3,3.0);

  % ---------- 高 A1O ----------
  \draw[dashed] (2.3,3.0) -- (2.3,0.6);

  % ---------- 顶点标注 ----------
  \node[below left]  at (0,0) {$A$};
  \node[below right] at (3,0) {$B$};
  \node[below right] at (4.62,1.12) {$C$};
  \node[left]        at (1.30,1.45) {$D$};
  \node[left]        at (2.16,3.04) {$A_1$};
  \node[above]       at (5.3,3.08) {$B_1$};
  \node[above right] at (6.92,4.22) {$C_1$};
  \node[above]       at (3.9,4.24) {$D_1$};
  \node[below]       at (2.12,0.42) {$O$};

\end{tikzpicture}
